% !TEX program = pdflatex
\documentclass[11pt,a4paper]{article}

% =========================================================================
% COMPARACIÓN FINANCIERA DE LA CARTERA
% Proyecto: CasaScore y Horus Sentinel
% Gestión de Proyectos TICS (EST315) -- Semestre 2026-II -- Team Sobrevivientes
% =========================================================================

% --- PAQUETES DE CONFIGURACIÓN ---
\usepackage[utf8]{inputenc}
\usepackage[T1]{fontenc}
\usepackage[spanish,es-tabla]{babel}
% Compatibilidad con babel-spanish reciente (flechas de TikZ y signo de porcentaje)
\AtBeginDocument{\shorthandoff{<>}\spanishplainpercent}
\usepackage{geometry}
\geometry{
    top=3.6cm,
    bottom=2.3cm,
    left=2.2cm,
    right=2.2cm,
    % PARTE MODIFICABLE: headheight debe ser mayor que la altura del logo más alto (2.56cm = 73pt).
    headheight=80pt,
    headsep=15pt,
    footskip=20pt
}

\usepackage{graphicx}
\usepackage{fancyhdr}
\usepackage[table]{xcolor}
\usepackage{titlesec}
\usepackage{booktabs}
\usepackage{tabularx}
\usepackage{array}
\usepackage{amsmath}
\usepackage{enumitem}
\usepackage{microtype}
\usepackage{amssymb}
% --- RECURSOS VISUALES: GRÁFICOS Y CÓDIGO DE COLORES ---
\usepackage{tikz}
\usetikzlibrary{positioning,calc}
\definecolor{GrisFila}{RGB}{246, 248, 251}
\definecolor{CuadCerca}{RGB}{150, 30, 30}
\definecolor{CuadSatisf}{RGB}{180, 130, 20}
\definecolor{CuadInfo}{RGB}{10, 45, 90}
\definecolor{CuadMonit}{RGB}{110, 115, 125}
\newcommand{\etiqueta}[2]{{\setlength{\fboxsep}{2pt}\colorbox{#1}{\textcolor{white}{\scriptsize\bfseries #2}}}}
\newcommand{\cerca}{\etiqueta{CuadCerca}{Gestionar de cerca}}
\newcommand{\satisf}{\etiqueta{CuadSatisf}{Mantener satisfecho}}
\newcommand{\informado}{\etiqueta{CuadInfo}{Mantener informado}}
\newcommand{\monit}{\etiqueta{CuadMonit}{Monitorear}}
\usepackage{needspace}
\usepackage{longtable}
\definecolor{Rama1}{RGB}{10, 45, 90}
\definecolor{Rama2}{RGB}{40, 95, 150}
\definecolor{Rama3}{RGB}{180, 130, 20}
\definecolor{Rama4}{RGB}{30, 110, 60}
\definecolor{Rama5}{RGB}{110, 60, 140}
\definecolor{RaciR}{RGB}{150, 30, 30}
\definecolor{RaciA}{RGB}{10, 45, 90}
\definecolor{RaciC}{RGB}{180, 130, 20}
\definecolor{RaciI}{RGB}{140, 145, 155}
\definecolor{NivAlto}{RGB}{30, 110, 60}
\definecolor{NivMedio}{RGB}{180, 130, 20}
\definecolor{NivBajo}{RGB}{150, 30, 30}
\usetikzlibrary{patterns}
\definecolor{SevBaja}{RGB}{30, 110, 60}
\definecolor{SevMedia}{RGB}{180, 130, 20}
\definecolor{SevAlta}{RGB}{150, 30, 30}
\usetikzlibrary{arrows.meta,shapes.geometric,fit,backgrounds}
\usepackage{listings}
\lstset{basicstyle=\ttfamily\scriptsize, frame=single, rulecolor=\color{AzulUNA!40}, backgroundcolor=\color{GrisFila}, columns=fullflexible, keepspaces=true, breaklines=true, xleftmargin=4pt, xrightmargin=4pt, literate={á}{{\'a}}1 {é}{{\'e}}1 {í}{{\'i}}1 {ó}{{\'o}}1 {ú}{{\'u}}1 {ñ}{{\~n}}1 {Á}{{\'A}}1}

\usepackage{hyperref}

% Logos: funciona a 2 o 3 niveles de profundidad dentro de docs/latex/
\graphicspath{{../assets/img/}{../../assets/img/}{../../../assets/img/}}

% --- PALETA INSTITUCIONAL FIEI / UNA PUNO ---
\definecolor{AzulUNA}{RGB}{10, 45, 90}
\definecolor{DoradoFIEI}{RGB}{180, 130, 20}
\definecolor{GrisOscuro}{RGB}{50, 50, 50}
\definecolor{AzulClaro}{RGB}{228, 236, 246}

% --- CONFIGURACIÓN DE LISTAS ---
\setlist{topsep=2pt, itemsep=2pt, parsep=1pt}

% --- COLUMNAS DE TABLA ---
\newcolumntype{L}{>{\raggedright\arraybackslash}X}
\newcolumntype{P}[1]{>{\raggedright\arraybackslash}p{#1}}
\newcolumntype{C}[1]{>{\centering\arraybackslash}p{#1}}
\newcolumntype{Q}[1]{>{\raggedleft\arraybackslash}p{#1}}
\newcolumntype{R}{>{\raggedleft\arraybackslash}X}
\renewcommand{\arraystretch}{1.25}

% --- FORMATO DE ENLACES ---
\hypersetup{
    colorlinks=true,
    linkcolor=AzulUNA,
    urlcolor=AzulUNA,
    citecolor=AzulUNA,
    pdftitle={Comparación Financiera De La Cartera -- CasaScore y Horus Sentinel},
    pdfauthor={Team Sobrevivientes}
}

% --- FORMATO DE TÍTULOS ---
\titleformat{\section}{\color{AzulUNA}\normalfont\large\bfseries}{\thesection.}{0.7em}{}[\color{AzulUNA}\titlerule]
\titleformat{\subsection}{\color{AzulUNA}\normalfont\normalsize\bfseries}{\thesubsection.}{0.5em}{}
\titlespacing*{\section}{0pt}{9pt}{4pt}
\titlespacing*{\subsection}{0pt}{6pt}{2pt}

% --- ENCABEZADO Y PIE DE PÁGINA (fancyhdr) ---
\pagestyle{fancy}
\fancyhf{}

% ===== TAMAÑO DE LOS LOGOS (PARTE MODIFICABLE) =====
% logo_fiei.png  : height=1.45cm (lado izquierdo)
% logo_equipo.png: height=2.56cm (lado derecho; 25 % más grande porque la imagen tiene más borde blanco)
% Para cambiar el tamaño, modifique solo el valor height=... de cada \includegraphics.
% Si el logo más alto supera 2.6cm, aumente también headheight y top en la configuración de página.
\fancyhead[L]{%
    \includegraphics[height=1.45cm,keepaspectratio]{logo_fiei.png}%
}
\fancyhead[C]{%
    \scriptsize \textbf{\textcolor{AzulUNA}{UNIVERSIDAD NACIONAL DEL ALTIPLANO DE PUNO}}\\[2pt]
    \tiny \textcolor{GrisOscuro}{FACULTAD DE INGENIERÍA ESTADÍSTICA E INFORMÁTICA}\\[2pt]
    \scriptsize \textbf{\textcolor{AzulUNA}{Gestión de Proyectos TICS (EST315) \textbar\ Team Sobrevivientes}}\\[2pt]
    \footnotesize \textbf{\textcolor{DoradoFIEI}{COMPARACIÓN FINANCIERA DE LA CARTERA}}%
}
\fancyhead[R]{%
    \includegraphics[height=2.56cm,keepaspectratio]{logo_equipo.png}%
}

\renewcommand{\headrulewidth}{0pt}
\renewcommand{\headrule}{}

\fancyfoot[L]{\scriptsize \textcolor{GrisOscuro}{CasaScore y Horus Sentinel}}
\fancyfoot[C]{\small \thepage}
\fancyfoot[R]{\scriptsize \textcolor{GrisOscuro}{Puno, Perú -- Semestre 2026-II}}
\renewcommand{\footrulewidth}{0.4pt}

\begin{document}

\begin{center}
    {\Large \textbf{\textcolor{AzulUNA}{COMPARACIÓN FINANCIERA DE LA CARTERA}}}\\[3pt]
    {\normalsize \textbf{CasaScore (AlquilaPuno) y Horus Sentinel}}\\[4pt]
    \footnotesize \textbf{Fecha:} lunes 21/09/2026 (semana 5) \quad \textbar\quad \textbf{Versión:} 1.0 \quad \textbar\quad \textbf{Fase:} Planificación
\end{center}

\vspace{0.1cm}

\section{Datos Generales}
\noindent\begin{tabularx}{\textwidth}{@{}P{4.2cm} L@{}}
\toprule
\textbf{Encargado por:} & Clyde Neil Paricahua Pari -- Project Manager \\
\textbf{Elaborado por:} & Ruth Karina Apaza Solis -- Analista de Planificación \\
\textbf{Propósito:} & Comparar con las mismas condiciones la rentabilidad de los dos proyectos de la cartera, como insumo para la decisión de continuidad del equipo como sponsor. \\
\textbf{Insumos:} & Viabilidad y Business Case, cronograma y presupuesto detallados de cada proyecto (semanas 3 y 4) \\
\bottomrule
\end{tabularx}

\section{Condiciones Comunes de la Comparación}
{\small
\rowcolors{2}{white}{GrisFila}
\noindent\begin{tabularx}{\textwidth}{@{}P{4.2cm} L@{}}
\toprule
\rowcolor{AzulClaro}
\textbf{Condición} & \textbf{Valor aplicado a ambos proyectos} \\
\midrule
Inversión (año 0) & Estimación de costos del semestre: personal, horas extra, equipos y gastos directos. \\
Horizonte & 3 años de operación (2027 a 2029). \\
Tasa de descuento & 10\,\% anual. \\
Tarifas & Sueldo de mercado / 192 h $\times$ 1.09 (EsSalud), las mismas 4 tarifas en ambos proyectos. \\
Flujo financiero & Ingresos del proyecto menos costos de operación. \\
Flujo económico & Flujo financiero más el beneficio que reciben los usuarios. \\
\bottomrule
\end{tabularx}
}

\section{Flujos de Caja}
Valores en soles. El flujo económico de CasaScore crece con los usuarios activos (500, 1,000 y 2,000 al mes); el de Horus Sentinel es constante con 1,000 usuarios y depende del convenio municipal de S/ 6,000 al año.

\vspace{0.1cm}
{\small
\rowcolors{2}{white}{GrisFila}
\noindent\begin{tabularx}{\textwidth}{@{}C{1.0cm} R R R R@{}}
\toprule
\rowcolor{AzulClaro}
& \multicolumn{2}{c}{\textbf{CasaScore}} & \multicolumn{2}{c}{\textbf{Horus Sentinel}} \\
\rowcolor{AzulClaro}
\textbf{Año} & \textbf{Financiero} & \textbf{Económico} & \textbf{Financiero} & \textbf{Económico} \\
\midrule
0 & $-$11,108.07 & $-$11,108.07 & $-$11,829.93 & $-$11,829.93 \\
1 & $-$2,869.46 & 11,704.54 & $-$182.75 & 11,817.25 \\
2 & $-$2,836.49 & 26,311.51 & $-$182.75 & 11,817.25 \\
3 & $-$1,352.48 & 56,943.52 & $-$182.75 & 11,817.25 \\
\bottomrule
\end{tabularx}
}

\Needspace{16\baselineskip}
\section{Indicadores de Rentabilidad}
{\small
\rowcolors{2}{white}{GrisFila}
\noindent\begin{tabularx}{\textwidth}{@{}L Q{2.8cm} Q{2.8cm} C{2.6cm}@{}}
\toprule
\rowcolor{AzulClaro}
\textbf{Indicador} & \textbf{CasaScore} & \textbf{Horus Sentinel} & \textbf{Más favorable} \\
\midrule
VAN económico & S/ +64,059.98 & S/ +17,557.81 & CasaScore \\
TIR económica & 166.5\,\% & 83.8\,\% & CasaScore \\
Payback simple & 0.95 años & 1.00 año & CasaScore \\
Payback descontado & 1.02 años & 1.11 años & CasaScore \\
Relación beneficio/costo & 6.77 & 2.48 & CasaScore \\
ROI a 3 años & 754.9\,\% & 199.7\,\% & CasaScore \\
VAN financiero & S/ $-$17,077.02 & S/ $-$12,284.40 & Horus Sentinel \\
\bottomrule
\end{tabularx}
}

\vspace{0.1cm}
\noindent Ninguno de los dos proyectos recupera la inversión con sus propios ingresos en 3 años (VAN financiero negativo); ambos se justifican por el beneficio que generan a sus usuarios. El VAN financiero de Horus Sentinel es menos negativo solo porque incluye los S/ 6,000 anuales del convenio municipal, que aún no está firmado.

\Needspace{24\baselineskip}
\section{Flujo Económico Descontado Acumulado}
\begin{center}
\begin{tikzpicture}[x=3.2cm,y=0.075cm,font=\scriptsize]
\foreach \v in {-20,0,20,40,60} {\draw[gray!25] (0,\v) -- (3,\v); \node[anchor=east] at (-0.05,\v) {\v};}
\draw[GrisOscuro, thick] (0,0) -- (3,0);
\foreach \a in {0,1,2,3} {\node[below] at (\a,-22) {Año \a};}
\node[rotate=90] at (-0.32,20) {Miles de S/};
\draw[AzulUNA, very thick] (0,-11.108) -- (1,-0.468) -- (2,21.277) -- (3,64.060);
\foreach \x/\y in {0/-11.108,1/-0.468,2/21.277,3/64.060} {\fill[AzulUNA] (\x,\y) circle (1.8pt);}
\draw[DoradoFIEI, very thick] (0,-11.830) -- (1,-1.087) -- (2,8.679) -- (3,17.558);
\foreach \x/\y in {0/-11.830,1/-1.087,2/8.679,3/17.558} {\fill[DoradoFIEI] (\x,\y) circle (1.8pt);}
\node[right, text=AzulUNA, font=\scriptsize\bfseries] at (3.02,64.06) {64.06};
\node[right, text=DoradoFIEI, font=\scriptsize\bfseries] at (3.02,17.56) {17.56};
\end{tikzpicture}
\\[2pt]
{\scriptsize \textcolor{AzulUNA}{\rule{12pt}{2pt}}~CasaScore \quad \textcolor{DoradoFIEI}{\rule{12pt}{2pt}}~Horus Sentinel \quad El punto en que cada línea cruza el cero es el payback descontado.}
\end{center}

\Needspace{16\baselineskip}
\section{Análisis de Sensibilidad}
VAN económico al 10\,\% cuando cambian los supuestos más inciertos de cada proyecto.

\vspace{0.1cm}
{\small
\rowcolors{2}{white}{GrisFila}
\noindent\begin{tabularx}{\textwidth}{@{}L Q{3.2cm} Q{3.2cm}@{}}
\toprule
\rowcolor{AzulClaro}
\textbf{Escenario} & \textbf{CasaScore} & \textbf{Horus Sentinel} \\
\midrule
Base & S/ +64,059.98 & S/ +17,557.81 \\
Beneficio para los usuarios 30\,\% menor & S/ +39,718.88 & S/ +8,605.15 \\
Sin convenio municipal & No aplica & S/ +2,636.71 \\
Sin convenio y beneficio 30\,\% menor & No aplica & S/ $-$6,315.96 \\
\bottomrule
\end{tabularx}
}

\vspace{0.1cm}
\noindent CasaScore se mantiene rentable aun con un beneficio 30\,\% menor (TIR de 115.7\,\%). Horus Sentinel depende del convenio: sin él, su VAN económico cae a S/ +2,636.71 con una TIR de 22.3\,\% y un payback de 2.03 años, y se vuelve negativo si además el beneficio es 30\,\% menor.

\Needspace{16\baselineskip}
\section{Factores de Ejecución en el Semestre}
{\small
\rowcolors{2}{white}{GrisFila}
\noindent\begin{tabularx}{\textwidth}{@{}L Q{3.2cm} Q{3.2cm}@{}}
\toprule
\rowcolor{AzulClaro}
\textbf{Factor} & \textbf{CasaScore} & \textbf{Horus Sentinel} \\
\midrule
Duración de la ruta crítica & 16 semanas & 19 semanas \\
Probabilidad de terminar en la semana 17 & 75\,\% & menos de 0.1\,\% \\
Línea base de costos & S/ 11,663.47 & S/ 12,421.43 \\
Presupuesto total con reserva de gestión & S/ 12,829.82 & S/ 13,663.57 \\
Gasto en efectivo por integrante & S/ 626.86 & S/ 807.33 \\
Dependencia externa & Ninguna & Convenio municipal \\
Origen de los datos & Levantamiento propio & Reportes ciudadanos \\
\bottomrule
\end{tabularx}
}

\Needspace{10\baselineskip}
\section{Lectura de los Resultados}
\begin{itemize}[leftmargin=*]
    \item CasaScore es más favorable en seis de los siete indicadores y genera 3.6 veces el VAN económico de Horus Sentinel con una inversión S/ 721.86 menor.
    \item La ventaja financiera de Horus Sentinel depende de un ingreso que no está asegurado.
    \item Horus Sentinel no termina dentro del semestre con la dedicación aprobada, por lo que su beneficio no llegaría a los usuarios en el plazo del curso.
    \item El resultado se presenta al equipo como sponsor junto con el informe de control de Horus Sentinel, para decidir su continuidad en la reunión de control del martes 22/09/2026.
\end{itemize}

\Needspace{22\baselineskip}
\section{Firmas}
\noindent En señal de conformidad con el presente documento, los integrantes del equipo firman a continuación:

\vspace{1.5cm}
\begin{center}
\begin{tabular}{c c c}
    \parbox{6.5cm}{\centering \rule{6.0cm}{0.6pt}\\[5pt] \textbf{Clyde Neil Paricahua Pari}\\[2pt] {\footnotesize Project Manager}}
    & \hspace{0.8cm} &
    \parbox{6.5cm}{\centering \rule{6.0cm}{0.6pt}\\[5pt] \textbf{Ruth Karina Apaza Solis}\\[2pt] {\footnotesize Analista de Planificación}}
    \\[2.0cm]
    \parbox{6.5cm}{\centering \rule{6.0cm}{0.6pt}\\[5pt] \textbf{Eliseo Tarqui Ajahuana}\\[2pt] {\footnotesize Especialista Técnico / Arquitectura}}
    & \hspace{0.8cm} &
    \parbox{6.5cm}{\centering \rule{6.0cm}{0.6pt}\\[5pt] \textbf{Viviana Mary Quisbert Quispe}\\[2pt] {\footnotesize Gestora de Comunicación, Riesgos y Documentación}}
\end{tabular}
\end{center}

\end{document}
